\section{WC 3/2/25}

\subsection*{Overview:}
This week I began coding my final artifact. Implementing my findings from my prior research and prototypes I feel very successful in my current progress

\subsection*{Achievements, Thoughts and Ideas:}
My first priority for this project was to build the authentication system for handling user login and management. I wanted to start here because I had already researched the principles behind it, and the UI was relatively simple—giving me a great opportunity to put my learnings into practice while ensuring the web app supports a wide range of devices. Having explore Shadcn previoulsy i used the components (alongside a lot of custom code) to produce the dynamic login page found here: \nameref{prototype:NewAuthPage}.

Once I had built the visual layer, I moved on to implementing a ‘Pinia' store. This powerful piece of technology acts as a global ‘database' within the browser, allowing every component in my code to check the user's state. This enables features such as a personalized welcome banner with the user's name, easy access to personalized pages, and—most importantly—a secure way to verify which user is accessing which data. The same system will also store users' graphs and notes globally, preventing unnecessary reloads when switching screens. Instead of repeatedly fetching data from the server, the app can use preloaded information, reducing server load and making it easier to support a larger number of users with fewer resources.

After completing the login page, I turned my attention to the sidebar. This serves as the central ‘spine' of my app, allowing users to navigate between pages and access key menus like settings. While many of the functions aren't fully operational yet, I successfully implemented account switching and logging out—both made possible by my earlier work on the backend Pinia store.

Next, I added a dark mode. Thanks to well-structured programming practices from earlier in development, this was relatively simple—I only needed to introduce a few additional components for theme switching and selected a couple of temporary colours to use (before I later work on branding and unifying the design)

In lessons this week, we worked on priority matrices, which helped me plan out which parts of my project needed to be completed first. Using this insight, I refined my project timeline into more detailed stages.


\subsection*{Power of Three:}
\textbf{This week's Tasks laid out per day}

\begin{tabularx}{\textwidth}{
  >{\raggedright\arraybackslash}X  % First column expands
  >{\centering\arraybackslash}p{1.5cm}  % Middle column has fixed width
  >{\raggedright\arraybackslash}X  % Third column expands
}
\toprule
\textbf{Goal} & \textbf{Achieved?} & \textbf{Notes} \\
\midrule


\rowcolor{gray!10}
\multicolumn{3}{l}{\textbf{Monday}} \\
Build Login Backup system (implementing prior prototype)  & \ding{51} & \\
Code Login UI & \ding{51} & \\
Brief Security Tests  & \ding{51} & Will conduct more stringent security audit before alpha release \\



\rowcolor{gray!10}
\multicolumn{3}{l}{\textbf{Tuesday}} \\
Implement Pinia store & \ding{51} & \\
Hook login system to it for application wide authentication & \ding{51} & \\
Begin coding sidebar, ensuring mobile responsibility & \ding{51} &  \\



\rowcolor{gray!10}
\multicolumn{3}{l}{\textbf{Thursday}} \\
Build Pinia tab store to store list of users tabs & \ding{51} & \\
Complete sidebar with ability to add and switch custom tabs in addition to the default tabs & \ding{51} & \\
Implement Dark mode & \ding{51} & Critical feature implemented\\  


\rowcolor{gray!10}
\multicolumn{3}{l}{\textbf{Thursday}} \\
Learn about prioritisation of tasks & \ding{51} & \\
Explore Priority matrixes & \ding{51} & \\
Update my Project choice page to include this matrix & \ding{51} &  Please find here: \nameref{projectchoice:prioritymatrix}\\

\bottomrule
\end{tabularx}


\subsection*{Plans for next week:}
In the next couple of weeks I hope to make more progress on the functionality of my webapp and make progress towards a fully functioning prototype In the near future
